% v6 section. Numbers are generated \vsix macros (E8, E8b, E2 evidence).
\section{What a Specification Costs, and What It Cannot Stop}\label{sec:v6cost}

The previous section shows that a gate enforces exactly what is written. On a public benchmark, that raises two
questions: what does writing it cost in legitimate work, and what happens to an attack that the writing did not
anticipate? Earlier versions of this paper answered the first with permissions copied from each task's reference
solution, which no deployment has. Here the permissions are written once per suite, without the solutions.

\para{Setup (E8, E8b).} AgentDojo v1.2.2~\cite{agentdojo}, banking (16 user tasks) and Slack (21), with AgentDojo's
``important instructions'' attack placed at each task's own injection vectors and AgentDojo's own graders. The
agent is \vsixEEightCkpt{} at temperature 0 in a JSON tool loop. Five arms see identical prompts and receive the same
reply, ``not permitted'', whenever a call is refused:
\begin{itemize}
  \item \emph{unprotected};
  \item \emph{Progent}, using its released code with a per-task policy generated by the same checkpoint~\cite{progent};
  \item \emph{strict}, a written specification \(S\) under which an effect may name an account, address, URL, user or
        channel only if it appears in the user's prompt or in the user's own state, and record writes must use values
        found there;
  \item \emph{delegated}, which is \(S\) plus trust in the output of any file or page the user's prompt names;
  \item \emph{oracle}, with tools, targets and amounts copied from the reference solution (a reference for what
        the finest specification admits, not deployable, and not an upper bound on utility).
\end{itemize}
Each user task runs once without attack and twice under assigned injections (\vsixEEightEpisodes{} episodes). E8b
repeats the attacked episodes with six injections we wrote to stay inside \(S\): pay an existing payee, raise an
existing standing order, schedule a recurring payment to an existing payee, message a secret to an existing user, add an
existing user to a private channel, and remove a user. Both plans,
the injections and their graders were committed before data.

\para{The cost before any model runs.} Offering each task's reference solution to each specification, with no model
involved, gives the specification-granularity cost directly. \(S\) admits \vsixStageStrictUser{} reference solutions
and \vsixStageStrictAttack{} of AgentDojo's reference attacks. Delegation admits \vsixStageDelegatedUser{} and
\vsixStageDelegatedAttack. The oracle admits \vsixStageOracleUser{} and \vsixStageOracleAttack. For E8b's injections
the same replay shows \(S\) admitting \vsixStageInsideStrict{} and the oracle \vsixStageInsideOracle.

\para{Live.} Figure~\ref{fig:v6spec}b plots benign utility against attack success. Strict \(S\) completed
\vsixStrictUtility{} benign tasks against \vsixUnprotectedUtility{} unprotected and \vsixOracleUtility{} with the
oracle grant; delegation completed \vsixDelegatedUtility. AgentDojo's injections succeeded in \vsixStrictNative{}
strict episodes, \vsixDelegatedNative{} delegated, \vsixOracleNative{} under the oracle grant and
\vsixUnprotectedNative{} unprotected. Injections inside \(S\) succeeded in \vsixStrictInside{} strict and \vsixDelegatedInside{} delegated
episodes, against \vsixUnprotectedInside{} unprotected, \vsixProgentInside{} under Progent and \vsixOracleInside{} under
the oracle grant. \vsixEEightProgentSentence

\para{Reading.} \vsixEEightCostSentence{} \vsixEEightPairedSentence{} \(S\) was written knowing AgentDojo's public tasks and injections, and the replay
without a model shows which reference attacks it refuses before any model runs, so its native attack success counts only attacks that
reach the injection's goal by a path its reference does not take. The informative numbers are therefore the utility
\(S\) costs and the attack success it leaves when the attacker stays inside what \(S\) permits. Those injections
succeeded in \vsixStrictInside{} strict episodes against \vsixUnprotectedInside{} unprotected; \(S\) refused none of
them, since it admits all of E8b's reference attacks (\vsixStageInsideStrict), so the difference is what the agent
did not carry out on its own. Progent's policies, generated per task by the same checkpoint from the user's request,
let \vsixProgentInside{} succeed, and the oracle, which admits only each solution's targets, \vsixOracleInside. On this
benchmark, then, our per-suite specification did no better than a per-task policy on benign tasks (the paired counts
above do not separate them) and let far more of the injections that stay inside it succeed. Those injections were written against \(S\), not
against Progent's policies, so this does not rank the two defences: it shows what an attack adapted to a specification
achieves, and an attack adapted to Progent's policies was not run. This is the paper's central hypothesis observed
against our own specification: a gate admits what its specification permits, and a specification written per task
permits less than one written per suite. AgentDojo has no concurrent calls, withdrawn grants or stops, so it
exercises none of what holding authority as state adds over a per-call policy, and nothing prevents enforcing a
per-task policy as state at the sink; what E8 and E8b show is the price of the specification's granularity, not of
where it is enforced.

\para{Across checkpoints (E2).} Re-running E8 unchanged with \vsixETwoCkpts{} as the agent, each served by the
same llama.cpp build on a second host's CPU (E8's \vsixEEightCkpt{} ran on the first), gives \vsixETwoSentence{} The
comparison is descriptive: the checkpoints differ in family, size, chat template and how well they follow the JSON
protocol, Progent's policies come from each checkpoint itself, and the host differs from E8's. Inference was not
equally repeatable across checkpoints: pairs of arms of one episode whose replies differed while everything the model
had received was identical numbered \vsixETwoDivergence{}, so not every difference between arms is caused by
the arm, and fewer of them are for the checkpoints with more such pairs. \vsixETwoQwenSentence{} Every arm is in
Table~\ref{tab:v6e2} and Figure~\ref{fig:v6e2} (Appendix~\ref{app:v6e2}).

\para{Limits.} Two of AgentDojo's four suites; one attack template; none of AgentDojo's own defences (tool filtering,
spotlighting, an injection detector) or the detection defences of Section~\ref{sec:gaps} as further arms; a repository-specific tool loop rather than
AgentDojo's native function calling; local open-weight models. \vsixEEightDivergenceSentence{} E8b's injections and graders are ours and were chosen to
lie inside \(S\), which the replay without a model confirms; E8b measures how often the agent carries them out, which
depends on the model as well as on \(S\), not how often attackers find such injections. The sink mediates only its own
records; AgentDojo's effects run in AgentDojo's in-memory environment, and the arm's predicate, not the sink, decides
whether a call runs.
